\subsection{例题：从表示到选择预测}

\begin{example}[CES 偏好的替代弹性]
令 $u(x)=\left[\alpha x_1^\rho+(1-\alpha)x_2^\rho\right]^{1/\rho}$。对任意给定效用水平，边际替代率为
$\mathrm{MRS}=\frac{\alpha}{1-\alpha}(x_1/x_2)^{\rho-1}$。当 $\rho<1$ 时，随 $x_1/x_2$ 上升，MRS 下降，体现凸偏好；参数 $\rho$ 只在给定函数形式下有解释，严格递增变换不会改变选择，但会改变边际效用数值。
\end{example}

\begin{proposition}[位似偏好的需求比例]
若效用函数一次齐次，价格与收入同时乘以 $t>0$，最优消费束不变；若只收入乘以 $t$，需求按比例扩大。
\end{proposition}
\begin{proof}
一次齐次性给出 $u(tx)=tu(x)$，因此同一条射线上的束具有按比例变化的效用。预算约束在价格与收入同时乘以 $t$ 时化为原约束乘以同一因子，约束集合不变，故选择不变。只收入扩大时，若 $x^*$ 可行，则 $tx^*$ 处于新预算集；齐次性与严格拟凹性保证最优比例保持。
\end{proof}

\begin{example}[效用表示与福利判断的区别]
两个消费者的选择都可由 $u_i(x)=x$ 表示，但把消费者甲效用提高 1 与把消费者乙效用提高 1 是否同等重要，不由表示定理决定。社会福利函数需要额外给出权重或分配原则；个体效用的序数性不能自动生成社会比较。
\end{example}

\begin{example}[连续性检验]
若 $u(x)=x_1+x_2$，上水平集 $\{x:u(x)\ge c\}$ 和下水平集 $\{x:u(x)\le c\}$ 都是闭半空间，故偏好连续。若定义 $u(x)=1$ 当 $x_1>0$、$u(x)=0$ 当 $x_1=0$，则在边界处发生跳跃，连续性失败；连续效用表示定理不能直接应用。
\end{example}

\begin{remark}
本章的例题应按“偏好公理—效用表示—函数性质—选择预测”的顺序阅读。先判断函数能否表示，再判断它是否单调、拟凹或位似，最后才讨论需求和福利，避免把解析形式直接当成经济含义。
\end{remark}

\subsection{同一例题的四步核验：表示、性质、需求与比较静态}
考虑两商品 CES 形式
\[
 u(x_1,x_2)=\left[\alpha x_1^\rho+(1-\alpha)x_2^\rho\right]^{1/\rho},
 \qquad 0<\alpha<1,\quad \rho<1,\ \rho\ne0.
\]
第一步只讨论表示：由于 CES 函数连续，且在 $x_i>0$ 时严格递增，
其上水平集和下水平集均可用极限序列检验闭性，因此它确实承载一个连续偏好。
这一结论并不依赖具体价格，也不等同于“需求已经被求出”。

第二步核验凸性。对固定的效用水平 $\bar u$，将约束写成
\[
 \alpha x_1^\rho+(1-\alpha)x_2^\rho=\bar u^\rho.
\]
隐式求导得到
\[
 \frac{dx_2}{dx_1}
 =-\frac{\alpha}{1-\alpha}
   \left(\frac{x_1}{x_2}\right)^{\rho-1}.
\]
当 $\rho<1$ 时，沿曲线向右移动，绝对斜率下降；因此曲线向原点凸，
对应的上水平集是凸集。这里必须注意：斜率下降是局部微分判断，
凸偏好还需要由整个上水平集的集合性质确认，不能只凭一段曲线外观下结论。

第三步求需求。拉格朗日函数为
\[
 \mathcal L=\left[\alpha x_1^\rho+(1-\alpha)x_2^\rho\right]^{1/\rho}
 +\lambda(m-p_1x_1-p_2x_2).
\]
一阶条件相除后得到
\[
 \frac{\alpha}{1-\alpha}\left(\frac{x_1}{x_2}\right)^{\rho-1}
 =\frac{p_1}{p_2},
\]
从而
\[
 \frac{x_1}{x_2}
 =\left[\frac{\alpha p_2}{(1-\alpha)p_1}\right]^{1/(1-\rho)}.
\]
将该比例与预算方程联立，即可得到两项 Marshallian 需求。这个推导显示，
价格比决定的是消费比例，收入决定的是共同尺度；后者正是位似性的需求含义。

第四步作反事实核验。若 $p_1$ 上升而 $p_2,m$ 不变，则上式中的 $x_1/x_2$
下降；若 $m$ 与两种价格同时乘以 $t>0$，比例不变且预算约束整体缩放，
需求束保持不变。由此可把“函数形式—偏好性质—选择预测—比较静态”四层
严格区分：前一层的结论不能未经证明直接替代后一层。

\begin{example}[线性与 CES 的边界比较]
取 $\rho\to1$，CES 的替代性增强，等值线趋向直线，最优解可能转为角点或整段；
取 $\rho\to0$，经适当归一化后趋向 Cobb--Douglas 形式，需求比例由参数和价格共同决定；
取 $\rho\to-\infty$，等值线趋向完全互补的折线，最优点由两个固定比例约束共同确定。
这些是同一函数族内部的边界比较，用来核对图形、符号和极限，不是另行引入新的理论主题。
\end{example}

\subsection{例题核验表：不要跳过中间的逻辑层}
对任意给定的效用表达式，建议按以下顺序核对。首先检查定义域：若函数含有分数次幂、对数或负指数，必须先说明 $x_i>0$ 还是允许边界点。其次检查表示关系：需要验证 $x\succeq y$ 当且仅当 $u(x)\ge u(y)$，而不是只验证几个数值例子。再次检查性质：单调性由偏导数或集合包含关系确认，凸偏好由上水平集确认，位似性则由无差异曲线沿射线的比例结构确认。最后才把价格和收入代入最优化问题。

以 Cobb--Douglas 形式 $u=x_1^\alpha x_2^{1-\alpha}$ 为例，定义域边界 $x_i=0$ 不能被简单略去。若 $0<\alpha<1$，内点偏导数存在且为正，但在边界处某一偏导可能发散；这不影响连续性，却会影响直接使用内点一阶条件的合法性。因此，完整证明必须分别说明：最优解为何存在、为何不是角点、为何一阶条件足以刻画全局最优。仅写出
\[
\frac{\alpha}{1-\alpha}\frac{x_2}{x_1}=\frac{p_1}{p_2}
\]
只能得到候选比例，不能独立构成需求定理。

\begin{proposition}[严格单调与财富耗尽的再核对]
若 $u$ 严格单调，预算集为非空紧集，且所有价格严格为正，则每个最优束均耗尽预算；若 $u$ 仅弱单调，结论可能失败。
\end{proposition}
\begin{proof}
严格单调情形可由增加一个商品的微小扰动证明。弱单调反例取 $u(x)=\min\{x_1,1\}$，当 $x_1\ge1$ 后增加商品一不再提高效用，因此存在未耗尽预算的最优束。这个反例说明“预算线上的最优”来自严格单调假设，而不是来自预算约束本身。
\end{proof}

\begin{example}[同一偏好下的不同数值表示]
若 $u(x)=x_1+x_2$，则 $v(x)=3u(x)+7$ 和 $w(x)=\log(1+e^{u(x)})$ 都保持相同的偏好排序。三者的水平集相同，但函数值差、边际效用和曲率不同。因此，看到某个表示的导数符号时，必须先问该结论是否在严格递增变换下保持；偏好顺序保持，而边际效用水平、效用差值和曲率通常不保持。
\end{example}
